% ========== MACRO ARCHITECTURE ==========
% Architectural layers, component boundaries, communication patterns, deployment topology
% Source: Symphony/Content/Architecture documents

\section*{Macro Architecture}
\addcontentsline{toc}{section}{Macro Architecture}

\lettrine{S}{ymphony's} macro architecture represents a paradigm shift from traditional monolithic IDE designs to a distributed, microkernel-inspired system that prioritizes modularity, performance, and extensibility. The architecture is built on four foundational principles: minimal core complexity, maximum extensibility, intelligent orchestration, and infrastructure-aware evolution.

\subsection*{Architectural Layers}
\addcontentsline{toc}{subsection}{Architectural Layers}

Symphony's architecture is organized into distinct layers, each with specific responsibilities and well-defined interfaces. This layered approach ensures separation of concerns while enabling efficient communication and resource sharing.

\subsubsection*{Presentation Layer}

The presentation layer encompasses all user-facing components and interfaces, built using modern web technologies for consistency across platforms.

\begin{compactlist}
    \item \textbf{React UI Framework}: Modern component-based interface with TypeScript integration
    \item \textbf{Jotai State Management}: Atomic state management for predictable UI behavior
    \item \textbf{TanStack Query}: Efficient data fetching and caching for responsive interactions
    \item \textbf{Tauri Integration}: Native desktop capabilities with web technology benefits
    \item \textbf{Communication Hub}: Centralized interface for backend communication
\end{compactlist}

\subsubsection*{Service Orchestration Layer}

The orchestration layer manages service coordination, task routing, and system-wide resource allocation through the Conductor microkernel.

\begin{compactlist}
    \item \textbf{Conductor Core}: Python-based RL model for intelligent decision-making
    \item \textbf{Task Router}: Distributes work across available services and extensions
    \item \textbf{Resource Manager}: Manages system resources and prevents conflicts
    \item \textbf{Conflict Arbitrator}: Resolves resource contention and access conflicts
    \item \textbf{Extension Registry}: Manages extension lifecycle and capabilities
\end{compactlist}

\subsubsection*{Infrastructure Layer}

The infrastructure layer provides core system services implemented as high-performance Rust extensions.

\begin{compactlist}
    \item \textbf{Pool Manager}: Manages AI model lifecycle and resource allocation
    \item \textbf{Artifact Store}: Content-addressable storage for intermediate results
    \item \textbf{DAG Tracker}: Workflow dependency tracking and execution coordination
    \item \textbf{Stale Manager}: Automated data lifecycle and archival management
    \item \textbf{Communication Bridge}: Cross-language interoperability infrastructure
\end{compactlist}

\subsubsection*{Extension Layer}

The extension layer hosts all specialized functionality as pluggable components, enabling infinite extensibility.

\begin{compactlist}
    \item \textbf{AI Model Services}: Specialized AI processing capabilities
    \item \textbf{Development Tools}: Language servers, debuggers, and analysis tools
    \item \textbf{Integration Services}: Version control, deployment, and external tool integration
    \item \textbf{Custom Extensions}: User-defined and community-contributed functionality
\end{compactlist}

\subsection*{Component Boundaries}

Symphony's component boundaries are designed to maximize modularity while minimizing coupling. Each component operates within well-defined interfaces and communication protocols.

\begin{infobox}[title=Boundary Design Principles]
\textbf{Interface Segregation}: Components expose only necessary functionality through minimal interfaces.

\textbf{Dependency Inversion}: High-level components depend on abstractions, not concrete implementations.

\textbf{Single Responsibility}: Each component has a single, well-defined purpose within the system.

\textbf{Open/Closed Principle}: Components are open for extension but closed for modification.
\end{infobox}

\subsubsection*{Core-Extension Boundary}

The boundary between Symphony's core and extensions is carefully designed to ensure system stability while enabling maximum flexibility.

\textbf{Core Responsibilities}:
\begin{compactlist}
    \item System initialization and shutdown
    \item Extension lifecycle management
    \item Resource allocation and conflict resolution
    \item Inter-component communication routing
    \item Security policy enforcement
\end{compactlist}

\textbf{Extension Responsibilities}:
\begin{compactlist}
    \item Specialized functionality implementation
    \item Domain-specific processing logic
    \item External system integration
    \item User interface customization
    \item Performance optimization
\end{compactlist}

\subsubsection*{Language Boundaries}

Symphony's multi-language architecture requires careful boundary management between Python, Rust, and JavaScript components.

\begin{table}[h]
\centering
\caption{Language Boundary Characteristics}
\label{tab:language-boundaries}
\begin{tabular}{@{}llll@{}}
\toprule
\textbf{Boundary} & \textbf{Mechanism} & \textbf{Performance} & \textbf{Use Case} \\
\midrule
Python-Rust & PyO3 FFI & 50-100ns overhead & AI model integration \\
Rust-JavaScript & Tauri IPC & 100-500ns overhead & UI communication \\
Python-JavaScript & WebSocket/HTTP & 1-5ms overhead & Real-time updates \\
\bottomrule
\end{tabular}
\end{table}

\subsection*{Communication Patterns}

Symphony employs multiple communication patterns optimized for different use cases and performance requirements.

\subsubsection*{Synchronous Communication}

Used for operations requiring immediate responses and strong consistency guarantees.

\textbf{Request-Response Pattern}:
\begin{compactlist}
    \item Direct function calls within language boundaries
    \item FFI calls for cross-language synchronous operations
    \item Blocking operations with timeout protection
    \item Error propagation and exception handling
\end{compactlist}

\subsubsection*{Asynchronous Communication}

Employed for operations that can be processed independently or in parallel.

\textbf{Event-Driven Pattern}:
\begin{compactlist}
    \item Publish-subscribe messaging for loose coupling
    \item Event queues for reliable message delivery
    \item Async/await patterns for non-blocking operations
    \item Callback mechanisms for completion notification
\end{compactlist}

\subsubsection*{Streaming Communication}

Utilized for real-time data flows and continuous processing scenarios.

\textbf{Stream Processing Pattern}:
\begin{compactlist}
    \item WebSocket connections for real-time UI updates
    \item Async iterators for data stream processing
    \item Backpressure handling for flow control
    \item Stream transformation and filtering
\end{compactlist}

\subsection*{Deployment Topology}

Symphony's deployment architecture supports multiple deployment scenarios while maintaining consistent behavior and performance characteristics.

\subsubsection*{Desktop Deployment}

The primary deployment target leveraging Tauri for native desktop integration.

\begin{successbox}
\textbf{Desktop Advantages}:
\begin{compactlist}
    \item Native file system access and performance
    \item Direct hardware resource utilization
    \item Offline operation capabilities
    \item Enhanced security through OS-level isolation
    \item Optimal AI model performance with local GPU access
\end{compactlist}
\end{successbox}

\subsubsection*{Web Deployment}

Browser-based deployment for accessibility and collaboration scenarios.

\textbf{Web Architecture Adaptations}:
\begin{compactlist}
    \item WebAssembly compilation for performance-critical components
    \item Service worker integration for offline capabilities
    \item Progressive Web App (PWA) features for native-like experience
    \item Cloud-based AI model execution for resource-intensive operations
\end{compactlist}

\subsubsection*{Hybrid Deployment}

Combination deployment supporting both local and cloud resources.

\textbf{Hybrid Capabilities}:
\begin{compactlist}
    \item Local processing for sensitive data and low-latency operations
    \item Cloud processing for resource-intensive AI model execution
    \item Intelligent workload distribution based on resource availability
    \item Seamless failover between local and cloud resources
\end{compactlist}

\begin{alertbox}
\textbf{Deployment Considerations}: Each deployment topology maintains the same architectural principles while adapting to platform-specific constraints and capabilities. The modular design ensures consistent behavior across all deployment scenarios.
\end{alertbox}